Un accidente no es provocado por una sola de estas fallas, sino por un conjunto de ellas, generalmente debidas a errores humanos. Dentro de estos, podemos decir que la falta de entrenamiento del técnico ha sido la causa de la mayoría de los accidentes, mientras que la negligencia ha provocado los más lamentables.

La improvisación en el desarrollo de las actividades, generalmente producto de la premura impuesta al personal o generada por la fatiga o la pereza y el exceso de confianza, son la causa más frecuente de las sobreexposiciones siempre innecesarias, tanto del personal como del público, sobre todo porque los efectos de radiación no se presentan inmediatamente después de haber recibido la exposición, por lo que las personas siguen trabajando de la misma manera, aumentando aún más la dosis que reciben.

El exceso de confianza también se refleja en accidentes de tráfico y pérdida de contenedores muchas veces de consecuencias lamentables.

La negligencia al no delimitar las áreas y/o no evacuarlas, es el motivo de la sobreexposición del público.

Sin embargo, es el desacoplamiento de las fuentes lo que genera el mayor número de accidentes y de las exposiciones más elevadas. Las consecuencias de un desacoplamiento dependen de las fallas que le acompañan. La falta o mal funcionamiento de los instrumentos de detección, es el segundo eslabón de la cadena y el tercero es la falta de verificación de que la fuente se encuentre en el contenedor. Esta última es la que ha propiciado las pérdidas de las fuentes, con las consecuencias más lamentables y las lesiones más graves a los técnicos.

En cuanto a la gravedad de las lesiones producidas, el aplastamiento de la manguera aunado a la ignorancia y la improvisación por la falta de equipo, constituye la segunda cadena más importante.

La descompostura de la chapa y el mal funcionamiento hasta ahora han sido sólo causas de sobreexposiciones del personal.

Como hemos mencionado, sólo la capacitación puede romper estas fallas, pues si bien es cierto que los accidentes se seguirán presentando, se reducirán las consecuencias de los mismos.
